\section{Crystallographic Structure and Interface Behaviour}

Interface orientation and underlying crystal symmetry strongly influence both the physical plausibility and ranking fidelity of heterostructures. Chapter~\ref{chapter:computational_investigation} demonstrates that coherent interfaces with well-aligned Miller planes tend to exhibit better MLP–DFT rank agreement, particularly in systems with low strain and high registry.

Future work will group interfaces by:
\begin{itemize}
    \item Parent symmetry (e.g., diamond vs. zincblende structures)
    \item Orientation combinations (e.g., \{100\}|\{100\}, \{100\}|\{111\})
    \item Slab asymmetry (e.g., Ge over Sn vs. Sn over Ge)
\end{itemize}

Metrics such as \( \rho \), \( \tau \), and Top-N overlap will be aggregated for each group to identify structural contexts that promote or degrade predictive performance. For example, early results suggest that silicon in the lower slab often reduces rank fidelity compared to reversed configurations.

If resources permit, this analysis could be extended to other lattice types (e.g., bcc or hcp) to evaluate whether the observed trends generalise beyond diamond-structured group 14 systems.